\subsection{ProtoPNet}
\label{subsubsec:eda_ppnet}

ProtoPNet (Prototypical Part Network) es una arquitectura de interpretabilidad \textit{ante-hoc} y \textit{model-specific}, ya que la explicación forma parte del propio mecanismo de decisión del modelo~\parencite{2019protopnet}. El mecanismo de decisión de este modelo consiste en comparar regiones de la representación interna de una muestra con un conjunto prototipos aprendidos durante el entrenamiento. Con esta idea, proporciona explicaciones locales de la forma ``esta parte de la entrada se asemeja a cierta región de este prototipo''.

Además de proporcionar explicaciones locales para muestras concretas, ProtoPNet también permite obtener una interpretación más global del modelo mediante el análisis de los prototipos asociados a cada clase. Al observar qué prototipos utiliza el modelo para reconocer una determinada categoría, es posible identificar qué patrones considera representativos de dicha clase. Al tratarse de un método de interpretabilidad basado en instancias (ejemplos reales del conjunto de datos), este enfoque guarda una estrecha relación metodológica con las explicaciones basadas en contraejemplos. Mientras que los prototipos de ProtoPNet justifican la decisión por similitud con casos típicos de la misma clase, los contraejemplos permiten entender los límites de decisión del modelo al mostrar qué variaciones mínimas en la instancia habrían alterado la clasificación hacia la categoría opuesta.

En el contexto de la detección de audios generados, una arquitectura inspirada en ProtoPNet podría aplicarse sobre representaciones tiempo-frecuencia, como espectrogramas, o sobre representaciones internas aprendidas por una red neuronal. En este caso, los prototipos podrían corresponder a patrones acústicos recurrentes, como determinadas estructuras espectrales, regiones temporales o zonas tiempo-frecuencia relevantes para distinguir entre audios \textit{bona fide} y \textit{spoof}. La explicación tendría entonces la forma: ``esta región del audio se parece a este patrón prototípico aprendido para una clase''.
